%*************************************************
% In this file the abstract is typeset.
% Make changes accordingly.
%*************************************************

\addcontentsline{toc}{section}{چکیده}
\newgeometry{left=2.5cm,right=3cm,top=3cm,bottom=2.5cm,includehead=false,headsep=1cm,footnotesep=.5cm}
\setcounter{page}{1}
\pagenumbering{arabic}						% شماره صفحات با عدد
\thispagestyle{empty}

~\vfill

\subsection*{چکیده}
\begin{small}
\baselineskip=0.7cm

با افزایش پیچیدگی تجهیزات صنعتی و هزینه‌های ناشی از توقف‌های پیش‌بینی‌نشده، توسعه سامانه‌هایی برای پایش پیوسته وضعیت، پیش‌بینی خرابی و اجرای تصمیم مناسب ضروری است. با وجود قابلیت دوقلوهای دیجیتال در اتصال دارایی‌های فیزیکی به بازنمایی‌های مجازی، یکپارچه‌سازی پیش‌بینی عمر مفید باقیمانده، پایش شبه‌بی‌درنگ، تعامل انسان و سامانه و تصمیم‌گیری ایمن در قالب معماری‌های توسعه‌پذیر همچنان یک چالش محسوب می‌شود. هدف این پژوهش، طراحی، پیاده‌سازی و ارزیابی یک چارچوب دوقلوی دیجیتال چندعاملی برای پایش و پیش‌بینی خرابی تجهیزات صنعتی با بهره‌گیری از یادگیری عمیق و مدل‌های زبانی بزرگ است. معماری پیشنهادی از پنج لایه موجودیت فیزیکی، گردآوری داده، یکپارچه‌سازی داده، عامل‌های دوقلوی دیجیتال و عامل‌های کاربردی تشکیل شده است. برای ارزیابی این چارچوب، دوقلوی دیجیتال موتورهای توربوفن بر پایه زیرمجموعه \lr{FD001} مجموعه‌داده
 \lr{CMAPSS}
ناسا توسعه یافته است. این زیرمجموعه شامل داده‌های ۱۰۰ موتور برای آموزش و ۱۰۰ موتور برای آزمون است. 
برای بررسی اثر تابع زیان بر رفتار
پیش‌بینی، دو پیکربندی از یک
معماری سه‌لایه مبتنی بر
\lr{LSTM} مورد ارزیابی قرار
گرفتند. مدل پایه با تابع
زیان \lr{MSE} به مقدار
\lr{RMSE} برابر ۱۵٫۹۳
و \lr{MAE} برابر ۱۱٫۹۵
دست یافت. در مدل دوم،
تابع زیان نامتقارن
\lr{Pinball} با ضریب
$\tau=0.125$ به کار
گرفته شد که مقدار
\lr{RMSE} برابر ۲۰٫۴۹
و \lr{MAE} برابر ۱۴٫۷۲
را ایجاد کرد. با وجود
افزایش خطای عددی،
نسبت کم‌برآورد در این
مدل از ۰٫۵۲ به
۰٫۷۱ افزایش یافت.
این نتیجه نشان می‌دهد
که تابع زیان نامتقارن
توانسته است جهت خطای
پیش‌بینی را به سمت
کم‌برآورد تغییر دهد؛
هرچند این تغییر با
افزایش خطای کلی همراه
بوده است.
عامل‌های مستقل پایش، پیش‌بینی، تابلو وضعیت و تعامل زبانی نیز برای به‌روزرسانی شبه‌بی‌درنگ وضعیت دوقلو، ارائه اطلاعات مورد نیاز و اجرای فرمان‌های عملیاتی از پیش‌ تعریف‌شده پیاده‌سازی شدند. نتایج نشان می‌دهد که چارچوب پیشنهادی، بستری دارای اجزای مستقل، قابل تعویض و توسعه‌پذیر برای نگهداری و تعمیرات پیش‌بینانه ایمن، تعاملی و قابل انطباق با دارایی‌ها و داده‌های صنعتی دیگر فراهم می‌کند.


\vspace*{0.7 cm}
\noindent\textbf{کلمات کلیدی:}
دوقلوی دیجیتال، سیستم چندعاملی، دوقلوهای دیجیتال چندعاملی، مدل زبانی بزرگ، نگهداری پیش‌بینانه، پیش‌بینی عمر مفید باقیمانده، 
\lr{LSTM}،
\lr{C-MAPSS}
ناسا،
 پایش بی‌درنگ، هوش مصنوعی توضیح‌پذیر.
\end{small}